\chapter{真空能量、暴胀与原始涨落}

\section{真空能量与德西特空间}

在平坦空间中，弗里德曼方程为
\begin{equation}
\left(\frac{\dot A}{A}\right)^2=\frac{8\pi G}{3}\rho.
\end{equation}
对于物质，\(\rho\)按\(A^{-3}\)的方式减小。对于辐射，\(\rho\)按\(A^{-4}\)的方式减小。对于真空能量，它并不稀释：
\begin{equation}
\rho=\rho_0=\text{constant}.
\end{equation}
如果宇宙中除了这种真空能量之外别无他物，则右侧为常数。将该常数写作\(H^2\)，
\begin{equation}
H^2=\frac{8\pi G}{3}\rho_0,
\qquad
H=\sqrt{\frac{8\pi G}{3}\rho_0}.
\end{equation}
于是哈勃常数真正是常数：
\begin{equation}
\frac{\dot A}{A}=H.
\end{equation}
因此
\begin{equation}
\dot A=HA,
\end{equation}
其解为指数形式：
\begin{equation}
A(t)=C e^{Ht}.
\end{equation}
相距遥远且不受束缚的星系之间的距离按指数增长。这就是德西特空间。

在这些坐标下，几何为
\begin{equation}
d\tau^2=dt^2-e^{2Ht}\left(dx^2+dy^2+dz^2\right).
\end{equation}
这里的真空能量为正。在平坦空间中它必须如此，因为弗里德曼方程左侧是非负的。负的真空能量只有在包含曲率项时才可能出现，而这里讨论的并非那种情形。

爱因斯坦引入宇宙学常数是出于另一个目的：平衡普通物质并使宇宙保持静止。

\section{当今加速膨胀的宇宙}

德西特空间也是当今宇宙在晚期的演化趋势。此前，标度因子遵循物质主导的行为
\begin{equation}
A(t)\propto t^{2/3}.
\end{equation}
几十亿年前，曲线向上弯曲并开始更快地增长。它还不是完美的指数，因为其他形式的能量仍然存在，但趋势是趋向于指数。

这里的加速意味着
\begin{equation}
\ddot A>0.
\end{equation}
这是晚期膨胀的一个观测特征。

现在考虑视界。哈勃定律说
\begin{equation}
v=HD.
\end{equation}
在退行速度变为光速的距离处，
\begin{equation}
D=\frac{c}{H}.
\end{equation}
等价地，可以将其改写为\(HD=c\)。\editorialnote{这是对口头关系\(D=c/H\)的编辑性代数重述。图~\ref{fig:lecture07-horizon}中右侧黑板上的上一行只有部分可辨认。}

\begin{figure}[ht]
\centering
\includegraphics[width=0.62\textwidth]{lecture_07_figure_02.png}
\caption{德西特膨胀中的宇宙视界关系\(D=c/H\)。}
\label{fig:lecture07-horizon}
\lectureframe{00:30:37}
\end{figure}

这个距离约为一百亿到一百五十亿光年，在数值上接近宇宙的年龄，只要将后者用\(c\)换算为距离。

\begin{classroomqa}
\lecturetimestamp{00:14:10}
\audiencequestion{视界尺度与宇宙年龄之间的近乎相等，难道不正是切线恰好经过原点时会发生的情况吗？}
\lecturerresponse{对于\(A(t)\propto t\)，那将完全成立。但一般情况下并不成立。在德西特空间中，\(H\)保持固定，而年龄不断增大，所以当前的近乎相等并非数学恒等式。}
\end{classroomqa}

\section{真空能量的含义}

这里用普通量子力学的语言来描述真空能量。对于谐振子，能量示意性地为
\begin{equation}
E\sim p^2+x^2.
\end{equation}
经典上，基态是平衡点，
\begin{equation}
x=0,\qquad p=0,\qquad E=0.
\end{equation}
量子力学不允许\(x\)和\(p\)同时精确地固定为零。不确定性原理迫使存在涨落，因此基态能量为正。对于谐振子，得到
\begin{equation}
E_0=\frac{1}{2}\hbar\omega.
\end{equation}

同样的论点随后被推广到场。宇宙充满了谐振子：电磁场的振荡，以及其他每一个场的振荡。每个场都有零点能。对于电磁场，能量密度示意性地为
\begin{equation}
u_{\mathrm{EM}}\sim E^2+B^2.
\end{equation}
即使不存在真实光子，仍然存在不可消除的零点振荡。这是对真空能量的一种解释。这里唯一需要的性质是，这种能量存在，并且不随宇宙膨胀而稀释。

\begin{classroomqa}
\lecturetimestamp{00:16:15}
\audiencequestion{这种真空能量是什么，为什么即使移除了物质和辐射，仍会有什么东西留存下来？}
\lecturerresponse{因为真空是量子场的基态，而量子基态不必具有零能量。谐振子具有正的基态能量，而场的行为如同振子的集合。所以即使不存在真实粒子，仍然存在零点能。}
\end{classroomqa}

这不是三度背景辐射。微波背景是普通辐射。它的能量密度减小，是因为光子稀释并且每个光子发生红移。真空能量则是在真实光子被移除之后仍然留存的残余零点能。

\section{\(H^{-1}\)意味着什么}

现在回到逆哈勃常数与宇宙年龄之间的关系。如果
\begin{equation}
A(t)\propto t,
\end{equation}
那么
\begin{equation}
H=\frac{\dot A}{A}=\frac{1}{t}.
\end{equation}
更一般地，如果
\begin{equation}
A(t)\propto t^p,
\end{equation}
那么
\begin{equation}
H=\frac{p}{t}.
\end{equation}
对于物质主导的情形 \(p=2/3\)，
\begin{equation}
H=\frac{2}{3t}.
\end{equation}
所以在幂律宇宙学中，\(H^{-1}\) 与年龄同量级，但并不恰好等于年龄。

德西特空间则不同。在那里
\begin{equation}
A(t)=C e^{\gamma t},
\end{equation}
因此
\begin{equation}
H=\frac{\dot A}{A}=\gamma,
\end{equation}
是一个真正的常数。这里 \(H^{-1}\) 并不是宇宙的年龄。它是指数增长的时间常数：即标度因子乘以 \(e\) 所需的时间。口头讨论中反复说 ``加倍''，随后又加以更正；精确的说法是乘以 \(e\)。今天，这个时间常数的量级为一百亿年。

\begin{classroomqa}
\lecturetimestamp{00:23:46}
\audiencequestion{您能否再讲一下关于宇宙年龄与逆哈勃常数的那一点？}
\lecturerresponse{对于幂律，可以得到 \(H=p/t\)，所以 \(H^{-1}\) 在量级为一的因子范围内与年龄相当。但在德西特空间中，\(H\) 是常数，而年龄却不断增长，因此 \(H^{-1}\) 只是指数时间常数，而不是年龄。}
\end{classroomqa}

\section{能量与宇宙视界}

如果真空能密度在宇宙膨胀时保持恒定，那么自然会问：真空总能量是否在增加？这里给出的回答是有限的。在广义相对论中，能量记账包含一个引力分量，因此总能量的概念比在初等力学中更为微妙。实际要点是与视界相关的那个。

在精确的德西特空间中，宇宙视界具有固定的大小，
\begin{equation}
D=\frac{c}{H}
=\frac{c}{\sqrt{8\pi G\rho_0/3}}.
\end{equation}
因此该视界内部的体积是固定的。其中的真空能为
\begin{equation}
E_{\mathrm{hor}}=\frac{4\pi}{3}D^3\rho_0.
\end{equation}
\editorialnote{这概括了 00：31：06--00：32：23 处的口头计算：取视界大小，求出球体体积，再乘以 \(\rho_0\)。}
由于 \(D\) 和 \(\rho_0\) 都是常数，这个量不随时间变化。对于某个给定观测者所能看到的宇宙部分而言，在精确的德西特空间中，可及的真空能保持固定。

事物也可以穿过宇宙视界，其意义与它们穿过黑洞视界大致相同。一旦到了外面，它们就与我们失去了因果联系。

\begin{classroomqa}
\lecturetimestamp{00:29:27}
\audiencequestion{如果真空能密度在宇宙膨胀时保持恒定，那么多出来的能量从何而来？}
\lecturerresponse{讲座用广义相对论的语言来回答，而不是用牛顿式的盒子语言。在精确的德西特空间中，视界大小保持固定，因此，在永远可被看到的区域内的真空能也保持固定。}
\end{classroomqa}

\section{为何引入暴胀}

随后，同样的德西特机制被推到遥远的过去，但时间常数要短得多，量级为
\begin{equation}
H^{-1}\sim 10^{-32}\,\mathrm{s}.
\end{equation}
这就是暴胀的思想：在任何直接可观测的事物之前，宇宙经历了一段指数膨胀时期。

基本动机是宇宙的平坦性和光滑性。空间切片看起来非常平坦。在大尺度上，宇宙也非常均匀。更引人注目的是，最后散射面当时就已经是光滑的，因此这种光滑化不可能只发生在后来可见的历史时期。

一个放瘪的气球是皱的。把它吹得足够大，任何可见的斑块都会变得光滑。葡萄干或西梅的类比说明了同一点：一个皱缩的物体，被拉伸一个巨大的因子后，在任何可见的斑块上都显得光滑。暴胀被认为正是做到了这一点。

这也以这里所用的形式回答了视界问题。最后散射面上相距遥远的区域对我们来说都是可见的，但在后来可见的历史时期，光没有时间从一个区域传播到另一个区域。因此，它们几乎均匀的密度和温度暗示着一段史前时期，其中一个较小的区域被极大地拉伸了。

暴胀有一种好坏参半的优点。它消除了对初始条件的敏感性，并且出于同样的原因，它也消除了关于宇宙如何起始的许多信息。

平坦性问题在这里以更尖锐的形式重新出现。从观测上看，曲率参数 \(k\) 看起来像零，精度大约在百分之一或更好。在可见区域内，几何看起来是欧几里得的。暴胀对此的解释是：使曲率半径远大于我们所能看到的区域。

\begin{classroomqa}
\lecturetimestamp{00:40:05}
\audiencequestion{我们是否假设暴胀之前处于平衡态？}
\lecturerresponse{不。暴胀的意义几乎相反。无论此前存在什么结构，都会被拉伸得如此之大，以至于可见区域失去了对更早细节的记忆。}
\end{classroomqa}

\section{标量场及其势}

为了解释早期宇宙的有效真空能如何可能远大于现在，引入一个标量场。与电场和磁场不同——它们是矢量——标量场没有方向。它只是在每一点有一个值：
\begin{equation}
\phi=\phi(x,y,z,t).
\end{equation}
这个场是假设性的。它并不等同于自然界中任何已知的场，其量子也不等同于任何已知的粒子。

在这里起作用的能量部分是势能，
\begin{equation}
V(\phi).
\end{equation}
所谓``势能''，指的是与场本身的值相关的能量，而不是与场的时间导数相关的能量。\(\phi\)的不同取值可以具有不同的能量。

\begin{figure}[ht]
\centering
\includegraphics[width=0.68\textwidth]{lecture_07_figure_03.png}
\caption{定性标量势\(V(\phi)\)，表明不同的场值携带不同的势能。}
\label{fig:lecture07-generic-potential}
\lectureframe{00:47:08}
\end{figure}

\begin{classroomqa}
\lecturetimestamp{00:46:34}
\audiencequestion{\(V(\phi)\)图中的横轴是什么？}
\lecturerresponse{横轴是场值\(\phi\)，纵轴是势能\(V\)。}
\end{classroomqa}

如果场在空间和时间上恒定，或在所关注的区域内几乎恒定，那么\(V(\phi)\)的行为就完全像真空能。它不会稀释。因此，先前的德西特公式可以复用，只需把\(\rho_0\)替换为\(V(\phi)\)。特别地，
\begin{equation}
H=\sqrt{\frac{8\pi G}{3}V(\phi)}.
\end{equation}

如果\(\phi\)随位置变化，不同区域会以不同的时间常数暴胀。但如果拉伸足够大，任何可见区域都会变得几乎均匀。

\begin{classroomqa}
\lecturetimestamp{01:05:11}
\audiencequestion{这里\(\phi\)是在充当时间参数吗？}
\lecturerresponse{在这一讨论中确实如此，其有限意义是：当场沿势滚下时，它是时间的单调函数。}
\end{classroomqa}

\section{平台、悬崖与再加热}

这一图景所需的势变得非常平坦，但并非完全平坦，然后到达悬崖边缘，急剧下降，并未完全降至零，随后又上升。

\begin{figure}[ht]
\centering
\includegraphics[width=0.68\textwidth]{lecture_07_figure_04.png}
\caption{具有高平台、陡降和较低极小值的暴胀势。}
\label{fig:lecture07-inflationary-potential}
\lectureframe{01:02:30}
\end{figure}

这种形状是高度人为构造的，但之所以引入它，是因为宇宙学图景似乎需要类似的东西。

场从平台高处开始。由于斜率很小，而与膨胀相关的有效黏滞很大，\(\phi\)变化得非常缓慢。一块缓慢穿过蜂蜜下落的石头，比一个在桌面上滑动的物体更接近所要描绘的图景。宇宙膨胀得越快，这种黏滞就越大。

\begin{classroomqa}
\lecturetimestamp{00:59:07}
\audiencequestion{是什么使\(\phi\)运动？}
\lecturerresponse{是势的斜率在推动它。斜率很小，而宇宙学黏滞很大，所以漂移是缓慢的，而不是迅速的。}
\end{classroomqa}

当\(\phi\)沿平台缓慢移动时，\(V(\phi)\)变化很小。因此，宇宙几乎就像真空能恒定那样暴胀。这一阶段的时间常数为
\begin{equation}
H^{-1}\sim 10^{-32}\,\mathrm{s}.
\end{equation}
一个粗略的口头估计把平台高度放在普朗克单位下的\(10^{-14}\)附近。\editorialnote{这是01:14:38--01:14:47处口头给出的``最佳猜测''估计，而不是讲座中推导出的值。}

暴胀持续了多久？这并不已知，因为平台的宽度和浅度并不已知。已知的是，宇宙至少被乘以\(e\)约五十次，而通常引用的数字是六十。\editorialnote{对01:14:05--01:19:27处口头讨论的e-折叠的编辑性重述。} 这些是e-折叠：每个e-折叠就是一次乘以\(e\)。

当场到达边缘时，势能迅速变化。场下落，其势能先转化为动能，然后转化为热和粒子，宇宙变得炽热。``再加热''这一名称在这一描述中略有误导，因为任何更早的热能都会在暴胀期间已经被稀释掉。

较低的那个极小值并不在零处。留在那里的那个小值，正是与当前缓慢的指数膨胀相关联的东西。

\begin{classroomqa}
\lecturetimestamp{01:19:48}
\audiencequestion{平台会不会恰好是平的，而扰动以某种方式把场从边缘推下去？}
\lecturerresponse{不会。它不可能完全平坦。它必须有一些倾斜。如果它太平坦，再加热后的宇宙会变得过于不均匀。}
\end{classroomqa}

``大爆炸''在这里用于两个相近的含义。它可以指再加热事件本身。它也可以指最早可直接看到的阶段，即最后散射面。稳定的要点是，实际观测到的一切都远低于暴胀平台。

\begin{classroomqa}
\lecturetimestamp{01:45:42}
\audiencequestion{再加热本身就是大爆炸吗？}
\lecturerresponse{如果你愿意，可以这么说。但那时宇宙仍然极其炽热且不透明。最后散射面要晚得多，是在宇宙冷却到足以变得透明之后。}
\end{classroomqa}

\section{量子涨落与星系的种子}

暴胀使可见宇宙非常均匀，但并非完全均匀。没有量子涨落，场会变得完全平滑。有了量子力学，它仍然会从一个地方到另一个地方略有变化。

中心图像把普通空间竖直画出，把场值\(\phi\)水平画出。如果场恰好是均匀的，那么在给定时刻，空间中每一点都会有相同的\(\phi\)，并且每一点都会一起向边缘漂移。于是每个区域都会在同一时刻从悬崖上掉下去。量子涨落破坏了这种精确的同步。

两幅图说明了这一点。一幅是在其他方面均匀的场中的一个小凸起。另一幅是一排士兵手挽手走向悬崖。如果他们完全对齐，他们就会一起掉下去。如果有些人稍微靠前、有些人稍微靠后，他们就会在不同时刻掉下去。

现在来了重要的效应。一个较早掉下去的区域会较早再加热。等到一个较晚的区域掉下去时，较早的那个区域已经有时间膨胀并稍微冷却。这种差异产生了再加热之后能量密度的变化。

\begin{classroomqa}
\lecturetimestamp{01:34:00}
\audiencequestion{较晚再加热的区域比更早的那些更热还是更冷？}
\lecturerresponse{更早的那些已经有时间膨胀并冷却。因此较晚的那些在掉下去时更热，而这种差异促成了密度不均匀性。}
\end{classroomqa}

观测到的效应很小：
\begin{equation}
\frac{\delta\rho}{\rho}\sim 10^{-5},
\end{equation}
大约十万分之一。这些是原初密度涨落，不是暗能量。这个数字在这里并不是仅凭定性草图推导出来的；它是用来限制平台能有多高、有多平坦的。

它们也是必需的。没有它们就不会有星系。再加热之后，引力放大了初始图案。过密区域从欠密区域吸引物质。过密区域变得更加过密，欠密区域变得更加欠密，其结果就是后来星系和恒星的结构。

这不仅仅是一个定性故事。人们可以计算膨胀宇宙中场的量子涨落，追踪它们经过再加热的过程，并把结果与从星系分布重建出的涨落图案进行比较。符合得非常好。

同样的讨论还引出一个关于熵的要点。引力成团看起来反直觉，但它并不违反第二定律。引力在统计力学中很特殊：它倾向于使东西成团，而这种成团提高熵，而不是降低熵。

\begin{classroomqa}
\lecturetimestamp{01:40:24}
\audiencequestion{既然\(\phi\)不是电磁场，微波背景在这一切中是如何体现的？}
\lecturerresponse{原初密度图案被印在最后散射面上。再加热把场能转化为普通粒子和光子，而微波背景后来携带了那些原初涨落的印记。}
\end{classroomqa}

\section{声波与下一步}

最后，讨论简要转向声波。这些声波出现在后来的引力成团之前。原初涨落被处理成声学振荡，而那些振荡在天空中留下具有已知物理尺度的团块。

这很重要，因为在已知距离上的已知物理尺度提供了一把量尺。一旦测量出它在最后散射面上所张的角度，就可以用几何来判断空间是平坦的、正曲率的还是负曲率的。细节留待以后。

\begin{classroomqa}
\lecturetimestamp{01:47:19}
\audiencequestion{声学波呢？}
\lecturerresponse{它们在天上提供了一把量尺。其已知的物理尺度，连同在最后散射面上观测到的角尺度，正是使人能够测量空间几何的东西。细节留待下一讲。}
\end{classroomqa}